% Section: A star lemma at cut-tight vertices
% Sources: problems/23/writeup/WALL_ATTACK_R41_GPTPRO56.md (sections 1-3);
%          problems/23/lean/Erdos23Delta0/Gamma/SingletonPairSigma.lean;
%          problems/23/lean/Erdos23Delta0/Gamma/CutTightStarPigeonhole.lean.
% Claims manifest: star_lemma_claims.md (same directory).

\section{A star lemma at cut-tight vertices}\label{sec:starlemma}

This section records a self-contained fact about maximum cuts of
triangle-free graphs. It uses only triangle-freeness and maximality of
the cut, and is independent of the support and Hall systems studied in
the rest of the paper.

Throughout the section, \(G\) is a finite simple graph with a fixed cut
\(V(G)=V_0\mathbin{\dot\cup}V_1\), \(B\) is the set of crossing edges,
and \(M=E(G)\setminus B\) is the set of monochromatic edges. The cut is
a \emph{maximum cut} when \(|M|\) is least over all cuts of \(G\), so
that \(|M|=\bip(G)\).

\begin{definition}[Switch loss; cut-tight vertices]\label{def:switchloss}
For \(S\subseteq V(G)\), let \(\partial_BS\) (respectively
\(\partial_MS\)) be the set of edges of \(B\) (respectively \(M\)) with
exactly one endpoint in \(S\), and define the \emph{switch loss}
\[
 \sigma(S)=|\partial_BS|-|\partial_MS|.
\]
For a vertex \(v\), write \(\sigma(v)=\sigma(\{v\})=d_B(v)-d_M(v)\),
where \(d_B(v)\) and \(d_M(v)\) count the crossing and monochromatic
edges at \(v\). We call \(v\) \emph{cut-tight} if \(\sigma(v)\le1\),
and we write \(N_B(v)\) for the set of neighbours joined to \(v\) by an
edge of \(B\), so \(|N_B(v)|=d_B(v)\).
\end{definition}

\begin{lemma}[Switch identity]\label{lem:switchloss}
Exchanging the shores of the vertices of a set \(S\subseteq V(G)\)
produces a cut with exactly \(|M|+\sigma(S)\) monochromatic edges.
Consequently, at a maximum cut,
\[
 \sigma(S)\ge0\quad\text{for every }S\subseteq V(G),
\]
and in particular \(\sigma(v)\ge0\) for every vertex \(v\).
\end{lemma}

\begin{proof}
An edge changes its crossing status under the switch if and only if
exactly one of its endpoints lies in \(S\); all other edges keep their
status. The new monochromatic set is therefore
\((M\setminus\partial_MS)\cup\partial_BS\), of size
\(|M|-|\partial_MS|+|\partial_BS|=|M|+\sigma(S)\). A maximum cut
minimizes the monochromatic count, so \(\sigma(S)\ge0\).
\end{proof}

\begin{lemma}[Additivity on a nonedge]\label{lem:nonedgeadd}
If \(x\ne y\) and \(xy\notin E(G)\), then for any cut
\[
 \sigma(\{x,y\})=\sigma(x)+\sigma(y).
\]
\end{lemma}

\begin{proof}
Since \(xy\notin E(G)\), no edge is incident to both \(x\) and \(y\).
Hence an edge has exactly one endpoint in \(\{x,y\}\) if and only if it
is incident to exactly one of \(x,y\), and
\(\partial_B\{x,y\}=\partial_B\{x\}\mathbin{\dot\cup}\partial_B\{y\}\),
\(\partial_M\{x,y\}=\partial_M\{x\}\mathbin{\dot\cup}\partial_M\{y\}\).
Subtract.
\end{proof}

\begin{lemma}[Star inequality]\label{lem:starineq}
Let \(G\) be triangle-free, fix any cut, and let \(v\) be a vertex with
\(k=d_B(v)\). Then
\[
 \sigma(\{v\}\cup N_B(v))
  =\sigma(v)+\sum_{a\in N_B(v)}\sigma(a)-2k,
\]
and hence, at a maximum cut,
\begin{equation}\label{eq:starineq}
 \sum_{a\in N_B(v)}\sigma(a)\;\ge\;2\,d_B(v)-\sigma(v).
\end{equation}
\end{lemma}

\begin{proof}
Set \(S=\{v\}\cup N_B(v)\). If two vertices \(a,a'\in N_B(v)\) were
adjacent, then \(va,va',aa'\) would form a triangle; so \(N_B(v)\) is
independent. An edge of \(G\) with both endpoints in \(S\) therefore
joins \(v\) to some \(a\in N_B(v)\), and since \(G\) is simple this
edge is the crossing edge \(va\) itself. Thus the edges of \(G\) inside
\(S\) are exactly the \(k\) crossing star edges, and no monochromatic
edge has both endpoints in \(S\). Summing degrees over \(S\) counts
each internal edge twice and each boundary edge once:
\[
 |\partial_BS|=\sum_{u\in S}d_B(u)-2k,
 \qquad
 |\partial_MS|=\sum_{u\in S}d_M(u).
\]
Subtracting gives the identity, and Lemma~\ref{lem:switchloss} gives
\eqref{eq:starineq} at a maximum cut.
\end{proof}

At a cut-tight vertex, \eqref{eq:starineq} reads
\(\sum_{a\in N_B(v)}\sigma(a)\ge2k-1\): the \(k\) crossing neighbours
carry average switch loss at least \(2-1/k\). Two immediate
consequences:

\begin{corollary}\label{cor:badendpoint}
At a maximum cut, every cut-tight vertex satisfies
\(d_M(v)\ge d_B(v)-1\). In particular, a cut-tight vertex with
\(d_B(v)\ge2\) is an endpoint of a monochromatic edge.
\end{corollary}

\begin{proof}
\(d_B(v)-d_M(v)=\sigma(v)\le1\).
\end{proof}

\begin{corollary}\label{cor:losstwoneighbour}
At a maximum cut of a triangle-free graph, every cut-tight vertex \(v\)
with \(d_B(v)\ge2\) has a crossing neighbour \(a\in N_B(v)\) with
\(\sigma(a)\ge2\).
\end{corollary}

\begin{proof}
Write \(k=d_B(v)\ge2\). If every \(a\in N_B(v)\) had \(\sigma(a)\le1\),
then \(\sum_{a\in N_B(v)}\sigma(a)\le k\le2k-2<2k-1\), contradicting
\eqref{eq:starineq} with \(\sigma(v)\le1\).
\end{proof}

A single high-loss neighbour already pairs with the whole star:

\begin{proposition}[Pair threshold]\label{prop:pairthreshold}
Let \(G\) be triangle-free with a maximum cut, and let \(x,y\) be
distinct crossing neighbours of a common vertex. If \(\sigma(x)\ge2\),
then \(xy\notin E(G)\) and
\[
 \sigma(\{x,y\})=\sigma(x)+\sigma(y)\ge2.
\]
\end{proposition}

\begin{proof}
Nonadjacency is triangle-freeness, as in the proof of
Lemma~\ref{lem:starineq}. By Lemma~\ref{lem:nonedgeadd},
\(\sigma(\{x,y\})=\sigma(x)+\sigma(y)\), and \(\sigma(y)\ge0\) by
Lemma~\ref{lem:switchloss}.
\end{proof}

The main statement refines Corollary~\ref{cor:losstwoneighbour}: the
combined loss of two neighbours can be forced \emph{relative to any
prescribed neighbour}. The arithmetic is a pigeonhole, which we isolate
in the form in which it has been formally verified.

\begin{lemma}[Star pigeonhole]\label{lem:starpigeon}
Let \(A\) be a finite set with \(|A|\ge2\), let
\(\ell\colon A\to\mathbb Z_{\ge0}\), and let \(s\in\{0,1\}\) satisfy
\[
 2|A|\;\le\;s+\sum_{a\in A}\ell(a).
\]
Then for every \(x\in A\) there exists \(y\in A\setminus\{x\}\) with
\(\ell(x)+\ell(y)\ge2\).
\end{lemma}

\begin{proof}
Suppose \(\ell(x)+\ell(y)\le1\) for every \(y\in A\setminus\{x\}\).
Since \(|A|\ge2\), some \(y_0\in A\setminus\{x\}\) exists; from
\(\ell(x)+\ell(y_0)\le1\) and \(\ell\ge0\) we get \(\ell(x)\le1\), and
likewise \(\ell(y)\le1\) for every \(y\in A\setminus\{x\}\). Hence
\(\sum_{a\in A}\ell(a)\le|A|\) and
\(2|A|\le s+|A|\le1+|A|\), so \(|A|\le1\), a contradiction.
\end{proof}

\begin{theorem}[Star lemma at cut-tight vertices]\label{thm:starlemma}
Let \(G\) be triangle-free, fix a maximum cut, and let \(v\) be a
cut-tight vertex with \(d_B(v)\ge2\). Then for every \(x\in N_B(v)\)
there exists \(y\in N_B(v)\setminus\{x\}\) with \(xy\notin E(G)\) and
\[
 \sigma(\{x,y\})=\sigma(x)+\sigma(y)\ge2.
\]
\end{theorem}

\begin{proof}
Apply Lemma~\ref{lem:starpigeon} with \(A=N_B(v)\), \(\ell=\sigma\)
restricted to \(A\) (nonnegative by Lemma~\ref{lem:switchloss}), and
\(s=\sigma(v)\in\{0,1\}\); hypothesis \eqref{eq:starineq} is the star
inequality. This yields \(y\in N_B(v)\setminus\{x\}\) with
\(\sigma(x)+\sigma(y)\ge2\). Since \(x\) and \(y\) are distinct
crossing neighbours of \(v\) and \(G\) is triangle-free, \(xy\notin
E(G)\), and Lemma~\ref{lem:nonedgeadd} gives
\(\sigma(\{x,y\})=\sigma(x)+\sigma(y)\).
\end{proof}

\begin{remark}
Theorem~\ref{thm:starlemma} is a necessary condition satisfied by every
maximum cut of every triangle-free graph. By itself it yields no upper
bound on \(\bip(G)\) and no progress on the Erd\H os \(N^2/25\)
conjecture; in this paper it serves only as a local production
threshold: inside the crossing star of any cut-tight vertex, every
neighbour can be paired to a nonadjacent partner whose joint switch
loss meets the threshold two.
\end{remark}

\begin{remark}[Formalization]\label{rem:leanstar}
Four ingredients of this section are verified in Lean~4, in the
development \texttt{Erdos23Delta0} accompanying this project. There,
graphs are literal edge lists with executable adjacency, cuts are side
assignments, and a maximum cut is a side assignment minimizing the
monochromatic count over all side assignments.
\begin{enumerate}[label=\textup{(\roman*)}]
\item Singleton nonnegativity (the vertex case of
Lemma~\ref{lem:switchloss}):
\texttt{singleton\_sigma\_nonneg\_of\_isMaxCut} in the module
\texttt{Erdos23Delta0.Gamma.SingletonPairSigma}; the general-set
inequality \(\sigma(S)\ge0\) is
\texttt{sigmaNonneg\_of\_badCount\_min} in
\texttt{Erdos23Delta0.CertGraph}.
\item Additivity on a nonedge (Lemma~\ref{lem:nonedgeadd}):
\texttt{sigma\_pair\_eq\_add\_singletons\_of\_nonadjacent} in
\texttt{Erdos23Delta0.Gamma.SingletonPairSigma}, with the crossing and
monochromatic boundary counts added separately
(\texttt{dB\_pair\_eq\_add\_singletons\_of\_nonadjacent},
\texttt{dM\_pair\_eq\_add\_singletons\_of\_nonadjacent}).
\item The pair threshold (Proposition~\ref{prop:pairthreshold}):
\texttt{nonadjacent\_of\_common\_blue},
\texttt{two\_le\_sigma\_pair\_of\_two\_le\_left}, and their combination
\texttt{common\_blue\_pair\_two\_le\_of\_left\_loss}, all in
\texttt{Erdos23Delta0.Gamma.SingletonPairSigma}.
\item The pigeonhole (Lemma~\ref{lem:starpigeon}):
\texttt{exists\_other\_with\_two\_le\_loss\_sum} in
\texttt{Erdos23Delta0.Gamma.CutTightStarPigeonhole}, stated over an
abstract finite neighbour set with \(\mathbb N\)-valued losses and the
star inequality as a hypothesis (the formal statement does not even
require the prescribed element to lie in the set).
\end{enumerate}
The geometric step, the star identity of Lemma~\ref{lem:starineq}, is
proved above but is not among the formalized statements; consequently
Theorem~\ref{thm:starlemma} as a single graph-level statement is proved
here on paper, with its arithmetic core and all pair manipulations
machine-checked.
\end{remark}
